% =====================================================================
%  Chapter 38 — The Verdict
% =====================================================================

\chapter{The Verdict, and What to Do With Your Own Money}
\label{ch:verdict}

The owner's question was whether there is a machine-learning procedure that
makes money and protects principal at the same time, and that can be
evaluated honestly enough that the investor knows which of the two it is
doing. This chapter answers it.

\section{The answer}

The system that this book builds earns \numFullCAGR{} a year against the
index's \numSpyCAGR{}, with a deepest drawdown of \numFullDD{} against
the index's \numSpyDD{}, a Sharpe ratio of \numFullSharpe{} against
\numSpySharpe{}, and a probability of backtest overfitting of \numPBO{}.

It does not beat the index, and the selection tests say clearly that it
could not be distinguished from the index on this evidence. It does
reduce the worst episode by a wide margin, and the risk governor that does
it contributes positively in every one of the \numWindowCount{} starting
months tested.

The machine-learning layer contributes \numTiltMean{} a year, with a
median of \numTiltMedian{}, and is positive in \numTiltPosFrac{} of the
windows.

So the answer to the question is, in three parts.

\textbf{Does it protect principal?} Not in the sense of never losing
money --- nothing does. In the sense that matters for a saver, which is a
bounded and recoverable worst episode: yes, substantially, and for a
reason that does not depend on a forecast being right.

\textbf{Does it make money?} It makes a real, positive compound return,
and it makes less than the index does. Whether that is ``making money''
depends on the reader's alternative, and the honest comparison is the
equal-weighted portfolio of the same universe, which makes more.

\textbf{Can it be evaluated honestly enough that the investor knows which
of the two it is doing?} Yes, and this is the part that was actually
novel. The book declares its trial family, publishes its count, applies
four selection tests to its own results, reports that they fail, and
computes its layer attribution from every starting month. A reader who
runs this can know exactly what they have.

\section{What the book actually claims}

Five claims, and they are not the ones the reader might expect.

\textbf{One: the prediction is the easy part and it does not work.} A
twelve-feature panel regression on price and volume, run over
\numModelConfigs{} configurations, produces a posterior over the forward
return whose signal-to-noise ratio is far below one. On this data, in
this universe, at this horizon, that is what a price-based forecast
looks like.

\textbf{Two: the risk layer is where the value is.} The governor needs no
forecast, so it cannot be wrong about the future, and its contribution is
stable across every window tested.

\textbf{Two and a half: the architecture is the point.} A layer that
prices its own uncertainty cannot destroy the portfolio it is part of.
The system's protection comes from the policy and the governor, and the
model's contribution to that protection is that it is bounded.

\textbf{Three: the best strategy in the universe was the naive one.} The
equal-weighted portfolio of sixty-one funds earned more compound return
over this record than the system, without a model, without an optimiser
and without a governor. The book's own evidence for this is in Chapter
25, from the literature, and in Chapter 32, from the record.

\textbf{Four: selection is the binding constraint on the whole subject.}
The probability of backtest overfitting above one half, and the implied
number of independent trials at \numImpliedTrials{} for a family of
\numFamilySize{}, are not statistics about this book. They are statistics
about the field.

\textbf{Five: the one large number in the subject that the reader can
change is behavioural.} The psychic cost of evaluating a long-horizon
portfolio annually is 5.1 per cent a year. No fee, no tax and no
execution error in this book is that large.

\begin{takeaway}[The book in one paragraph]
A machine-learning forecast of liquid equity ETFs, built honestly and
evaluated honestly, does not beat the index and cannot be distinguished
from it. A risk governor that needs no forecast reduces the worst episode
by seventeen percentage points and does so in every window tested. The
architecture that lets the two coexist --- a stated allocation the
machine cannot touch, a bounded tilt, a posterior that prices its own
uncertainty, a governor, a tax engine, and a trial family declared in
advance --- is the transferable result. The largest single number in the
subject is not a fee or a tax but the cost of looking at the account too
often.
\end{takeaway}

\section{What to do with your own money}

The book's recommendations, in order, and the reasons for them.

\textbf{Decide your own horizon before you decide your allocation.} The
strategy in this book is for a reader with twenty or more years and a
stable income. For a reader with five years, the tax and the behavioural
friction compound and the answer is probably the index, held.

\textbf{Write down the three numbers from Chapter 2.} The maximum
drawdown you will not abandon the strategy for, the number of years you
can wait for a recovery, and the annual withdrawal your plan depends on.
Those three numbers define the feasible set of portfolios, and they are
inputs to the system rather than outputs of it.

\textbf{If your circumstances permit it, use the tax-exempt account.} The
tax costs \numTaxPenalty{} a year of compound growth over this record.

\textbf{Do not expect the model to help you, and do not let it cost you.}
The tilt's contribution is \numTiltMean{} a year. If the reader runs the
system, the model's contribution is a rounding error, and the entire
value of running it is that it makes the discipline automatic.

\textbf{Look at the account at least four times a year, and not more.} This
is the one recommendation in the book with a number attached to it that is
larger than every other number in the book combined.

\textbf{Do not override it because the portfolio is down.} Override it when
your circumstances change. The strategic allocation was never the
model's to set, and it is the reader's to change.

\begin{takeaway}[The five sentences]
Write down your three numbers. Use the tax-exempt account if you can.
Do not expect the model to help and do not let it cost you. Look at the
account four times a year, not more. Do not sell in a drawdown.
\end{takeaway}

\section{The last thing this book has to say}

There is a version of this book in which the machine finds the edge and
the reader's problem is solved. It was not written, because the
measurements did not support it, and because a book that reports a null
result with the same rigour as an alpha is more useful to the reader than
one that reports only alphas.

But the null result is not nothing, and the reader should not read it as
an instruction to do nothing. The measurements say something specific:
that the protection layer works, that the prediction layer does not, and
that the second fact is not a reason to abandon the first. A portfolio
that is the stated policy plus a governor, with a bounded and honestly
uncertain forecast on top, gives a reader the index's risk-adjusted
return with a materially shallower worst episode, and it does so without
requiring the reader to believe anything that the evidence does not
support.

That is the patient part. It is also, on this evidence, all of the
available part.

\begin{keytakeaways}
\item The system protects the principal in the sense that matters --- a
bounded, recoverable worst episode --- and it does so without needing a
forecast to be right.
\item It makes less money than the index, and the best strategy in the
universe on this record was the naive equal-weighted one.
\item The evaluation is the transferable result: a declared trial family,
four selection tests applied to the book's own results, a layer
attribution from every starting month, and a null reported as a null.
\item The largest single number in the subject is behavioural: the cost of
evaluating the portfolio too often.
\end{keytakeaways}
